\chapter{Simple exponential smoothing}\label{ch:es}

This note documents \code{nabia::forecast::SimpleExponentialSmoothing}. We
define the recursion and prove that it is an exponentially weighted average whose
weights sum to one. We compute its mean lag and derive the span
convention $\alpha=2/(N+1)$. We then give the statistical model behind it, the
innovations state-space model ETS(A,N,N), and derive the forecast
distribution, including the variance $\sigma^2(1+(h-1)\alpha^2)$. We also prove
its equivalence with ARIMA(0,1,1) and with the steady-state Kalman filter of
the local level model. Finally, we describe the implementation: initialization,
variance estimate, prediction intervals and complexity.




%==============================================================================
\section{The recursion}

Let $y_0,\dots,y_{T-1}$ be an equally spaced series and $\alpha\in(0,1]$ the
smoothing factor. Simple exponential smoothing (SES)~\cite{brown1959} maintains a
\emph{level} $\ell_t$, updated as
\begin{equation}\label{es:eq:ses}
  \ell_t = \alpha\, y_t + (1-\alpha)\,\ell_{t-1},\qquad t=0,\dots,T-1,
\end{equation}
from an initial level $\ell_{-1}$. Equivalently, in \emph{error-correction form},
\begin{equation}\label{es:eq:error-correction}
  e_t = y_t - \ell_{t-1},\qquad \ell_t = \ell_{t-1} + \alpha\, e_t ,
\end{equation}
where $e_t$ is the one-step-ahead forecast error: $\ell_{t-1}$ is the forecast
of $y_t$ made at time $t-1$. Every forecast made at the end of the sample
equals the final level,
\[
  \hat y_{T-1+h\mid T-1} = \ell_{T-1},\qquad h=1,2,\dots
\]
SES is therefore suited to series without trend or seasonality: it estimates the
current level and extrapolates it flat. The value of $\alpha$ sets the trade-off
between reacting to recent data ($\alpha\to 1$, the naive forecast $\ell_t=y_t$)
and smoothing out noise ($\alpha\to 0$, an almost constant level).

\subsection{Exponential weights}

\begin{proposition}[Exponentially weighted average]\label{es:prop:weights}
For $t\ge 0$,
\begin{equation}\label{es:eq:weights}
  \ell_t = \sum_{j=0}^{t} \alpha(1-\alpha)^j\, y_{t-j} + (1-\alpha)^{t+1}\,\ell_{-1},
\end{equation}
and the weights add up to one:
$\sum_{j=0}^{t}\alpha(1-\alpha)^j + (1-\alpha)^{t+1} = 1$.
\end{proposition}
\begin{proof}
By induction on $t$. For $t=0$, \eqref{es:eq:ses} gives $\ell_0=\alpha y_0+(1-\alpha)\ell_{-1}$,
which is \eqref{es:eq:weights}. If \eqref{es:eq:weights} holds for $t-1$, then
\[
  \ell_t = \alpha y_t + (1-\alpha)\Big[\sum_{j=0}^{t-1}\alpha(1-\alpha)^j y_{t-1-j} + (1-\alpha)^t\ell_{-1}\Big]
         = \alpha y_t + \sum_{j=1}^{t}\alpha(1-\alpha)^{j} y_{t-j} + (1-\alpha)^{t+1}\ell_{-1},
\]
which is \eqref{es:eq:weights} for $t$. For the sum, the geometric series gives
$\sum_{j=0}^{t}\alpha(1-\alpha)^j = \alpha\,\frac{1-(1-\alpha)^{t+1}}{1-(1-\alpha)} = 1-(1-\alpha)^{t+1}$.
\end{proof}

So the level is a weighted average of past observations, with weights that decay
geometrically with age. The weight on the initial level $\ell_{-1}$ vanishes as $t$ grows, so
the choice of initialization matters only for the first $O(1/\alpha)$ steps.

\begin{proposition}[Mean lag and the span convention]\label{es:prop:lag}
In the limit $t\to\infty$, the weights $w_j=\alpha(1-\alpha)^j$ have mean age
$\bar{\jmath} = \sum_j j\,w_j = (1-\alpha)/\alpha$. A simple moving average over
the last $N$ points has mean age $(N-1)/2$. The two match when
$\alpha = 2/(N+1)$.
\end{proposition}
\begin{proof}
Differentiating the geometric series $\sum_{j\ge0}q^j=1/(1-q)$ gives
$\sum_{j\ge0} j q^{j-1}=1/(1-q)^2$. With $q=1-\alpha$,
$\sum_j j\,\alpha q^j = \alpha q/(1-q)^2 = (1-\alpha)/\alpha$. The moving average
puts weight $1/N$ on ages $0,\dots,N-1$, so its mean age is $(N-1)/2$. Setting
$(1-\alpha)/\alpha=(N-1)/2$ and solving gives $\alpha=2/(N+1)$.
\end{proof}

This is the convention behind \code{EWMA::alpha\_from\_span} in
\code{nabia/signal}. The mean lag is also the delay with which the level
follows a linear trend: if $y_t = a + bt$, the steady-state level is
$\ell_t = a + b(t-\bar{\jmath})$. This lag is why trend strategies built on SES
or EWMA signals react late, and why trend models (local linear trend in the
Kalman filter note) extrapolate the slope instead.

%==============================================================================
\section{The statistical model: ETS(A,N,N)}

On its own, \eqref{es:eq:ses} is only an algorithm. To attach uncertainty to the
forecasts we need a statistical model whose optimal one-step predictor is SES.
The standard choice is the innovations state-space model with additive error,
no trend and no seasonality, ETS(A,N,N)~\cite{hyndman2008}:
\begin{equation}\label{es:eq:ets}
  y_t = \ell_{t-1} + \varepsilon_t,\qquad
  \ell_t = \ell_{t-1} + \alpha\,\varepsilon_t,\qquad
  \varepsilon_t \overset{\text{iid}}{\sim} \N(0,\sigma^2).
\end{equation}
A single error source drives both equations (hence ``innovations'' form).
Given $\ell_{t-1}$, the conditional mean of $y_t$ is $\ell_{t-1}$, and
$\varepsilon_t = y_t-\ell_{t-1} = e_t$. Substituting into the state equation gives
exactly \eqref{es:eq:error-correction}. So the SES levels are the model's
states, and the SES errors are its innovations.

\begin{theorem}[Forecast distribution]\label{es:thm:forecast}
Under \eqref{es:eq:ets}, conditionally on the information up to time $T-1$
(that is, on $\ell_{T-1}$), for every horizon $h\ge 1$,
\[
  y_{T-1+h} \sim \N\!\big(\ell_{T-1},\; \sigma^2\,[1+(h-1)\alpha^2]\big).
\]
\end{theorem}
\begin{proof}
Iterating the state equation from $T-1$,
$\ell_{T-2+h} = \ell_{T-1} + \alpha\sum_{j=1}^{h-1}\varepsilon_{T-1+j}$. Then
\[
  y_{T-1+h} = \ell_{T-2+h} + \varepsilon_{T-1+h}
            = \ell_{T-1} + \alpha\sum_{j=1}^{h-1}\varepsilon_{T-1+j} + \varepsilon_{T-1+h}.
\]
The future innovations are independent of $\ell_{T-1}$ and of each other, with
mean $0$ and variance $\sigma^2$. The right-hand side is therefore Gaussian with mean
$\ell_{T-1}$ and variance $\alpha^2(h-1)\sigma^2+\sigma^2$.
\end{proof}

The point forecast is flat, but the uncertainty grows linearly in $h$. Each
future shock moves the level permanently by $\alpha\varepsilon$, so the model is a random walk observed with noise.
The parameter $\alpha$ sets how much of each shock is
permanent ($\alpha\varepsilon_t$ enters the level) and how much is transitory.

\subsection{Equivalence with ARIMA(0,1,1)}

\begin{proposition}\label{es:prop:arima}
Under \eqref{es:eq:ets}, the differences $\Delta y_t = y_t - y_{t-1}$ follow the
MA(1) process $\Delta y_t = \varepsilon_t - (1-\alpha)\,\varepsilon_{t-1}$, so $y$
is ARIMA(0,1,1) with moving-average coefficient $\theta=-(1-\alpha)$.
\end{proposition}
\begin{proof}
From \eqref{es:eq:ets},
$y_t - y_{t-1} = (\ell_{t-1}+\varepsilon_t)-(\ell_{t-2}+\varepsilon_{t-1}) = (\ell_{t-1}-\ell_{t-2}) + \varepsilon_t-\varepsilon_{t-1}
= \alpha\varepsilon_{t-1}+\varepsilon_t-\varepsilon_{t-1}$.
\end{proof}

The MA polynomial $1-(1-\alpha)z$ has its root at $1/(1-\alpha)$, outside the
unit circle for $\alpha\in(0,2)$. The process is therefore invertible, and $\varepsilon_t$ can be
recovered from past $y$. This is what makes the recursion \eqref{es:eq:ses}
forget its initialization (Proposition~\ref{es:prop:weights}). \code{nabia}
restricts $\alpha$ to $(0,1]$, where all weights in \eqref{es:eq:weights} are
non-negative and the level is a genuine average.

\subsection{Equivalence with the Kalman filter of the local level model}

The local level (random walk plus noise) model has two independent noise sources:
\begin{equation}\label{es:eq:local-level}
  x_t = x_{t-1} + w_t,\quad w_t\sim\N(0,q),\qquad
  y_t = x_t + v_t,\quad v_t\sim\N(0,r).
\end{equation}

\begin{proposition}[Muth~\cite{muth1960}]\label{es:prop:muth}
The Kalman filter for \eqref{es:eq:local-level} converges to the recursion
$\hat x_t = \hat x_{t-1} + K(y_t-\hat x_{t-1})$ with steady-state gain
\[
  K = \frac{P}{P+r},\qquad P = \frac{q+\sqrt{q^2+4qr}}{2}.
\]
This is SES with $\alpha=K$. Conversely, the SES with factor $\alpha\in(0,1)$ is the steady state of
\eqref{es:eq:local-level} with signal-to-noise ratio $q/r = \alpha^2/(1-\alpha)$.
\end{proposition}
\begin{proof}
For this model the predicted variance $P_t=\Var(x_t\mid y_{0:t-1})$ obeys the
Riccati recursion $P_{t+1} = P_t r/(P_t+r) + q$ (see the Kalman filter note,
Section~6). A fixed point satisfies $P(P+r)=Pr+q(P+r)$, that is
$P^2-qP-qr=0$, whose positive root is the $P$ above. The map
$P\mapsto Pr/(P+r)+q$ is increasing and concave with slope $r^2/(P+r)^2<1$ at the
fixed point, so the iteration converges to it from any $P_0\ge0$. The filtered
estimate is $\hat x_t = \hat x_{t-1} + K_t(y_t-\hat x_{t-1})$ with
$K_t=P_t/(P_t+r)\to K$. For the converse, $K=P/(P+r)$ gives $P=Kr/(1-K)$.
Substituting into $P^2=q(P+r)$ gives $K^2r^2/(1-K)^2 = q\,r/(1-K)$, that is $q/r=K^2/(1-K)$.
\end{proof}

The \code{nabia} test suite checks Proposition~\ref{es:prop:muth} numerically. After
300 steps, the local-level Kalman filter and an SES with $\alpha=K$ agree to
$10^{-9}$.

%==============================================================================
\section{Estimation}

\paragraph{Innovation variance.}
Conditionally on the initial level, the likelihood of \eqref{es:eq:ets} factorizes
over the one-step errors (prediction error decomposition):
$\log L(\sigma^2) = -\frac{T}{2}\log(2\pi\sigma^2) - \frac{1}{2\sigma^2}\sum_t e_t^2$.
Setting the derivative in $\sigma^2$ to zero gives the maximum-likelihood
estimate
\[
  \hat\sigma^2 = \frac{\mathrm{SSE}}{T},\qquad \mathrm{SSE}=\sum_{t=0}^{T-1} e_t^2 .
\]

\paragraph{Smoothing factor.}
For fixed $\sigma^2$, maximizing the likelihood in $\alpha$ amounts to minimizing
$\mathrm{SSE}(\alpha)$. This one-dimensional function is smooth but not
necessarily convex. \code{nabia} leaves the choice to the user: the
\code{forecast\_demo} example scans a grid, and \code{NelderMead} on $[0,1]$
or a golden-section search would also do.

%==============================================================================
\section{Implementation}

\subsection{Interface}
The class implements the forecaster interface
\[
  \code{IForecaster<SimpleExponentialSmoothingParams>}:
\]
\begin{lstlisting}
SimpleExponentialSmoothing model(0.3);           // or (std::make_unique<Params>(...))
model.fit(y);                                    // Eigen::VectorXd, finite values
Forecast fc = model.forecast(12);                // mean and variance per step
Eigen::VectorXd lo = fc.lower(0.95), hi = fc.upper(0.95);
Eigen::VectorXd one_step = model.fitted_values(); // l_{t-1}, the forecast of y_t
double sse = model.sse();
\end{lstlisting}

\subsection{Algorithm}
\begin{algorithm}[ht]
\caption{SES fit and forecast}\label{es:alg:ses}
\begin{algorithmic}[1]
\Require $y_0,\dots,y_{T-1}$ finite, $\alpha\in(0,1]$
\State $\ell\gets y_0$;\; $\mathrm{SSE}\gets0$ \Comment{$\ell_{-1}=y_0$}
\For{$t=0,\dots,T-1$}
  \State $\hat y_t\gets\ell$ \Comment{one-step forecast, stored in \code{fitted\_values()}}
  \State $e\gets y_t-\ell$;\; $\mathrm{SSE}\gets\mathrm{SSE}+e^2$;\; $\ell\gets\ell+\alpha e$
\EndFor
\State $\hat\sigma^2\gets\mathrm{SSE}/T$
\For{$h=1,\dots,H$} \Comment{\code{forecast(H)}}
  \State $\mu_h\gets\ell$;\; $v_h\gets\hat\sigma^2\,(1+(h-1)\alpha^2)$
\EndFor
\end{algorithmic}
\end{algorithm}

\paragraph{Initialization.}
The implementation uses $\ell_{-1}=y_0$, so the first error is $e_0=0$ and the
first fitted value equals the first observation. By
Proposition~\ref{es:prop:weights} the effect of this choice decays like
$(1-\alpha)^{t}$. Because $e_0=0$ by construction, $\hat\sigma^2=\mathrm{SSE}/T$
averages $T-1$ informative errors over $T$ terms. This slightly
underestimates $\sigma^2$ in short samples, a bias of factor $(T-1)/T$.

\paragraph{Prediction intervals.}
\code{Forecast::lower(level)} and \code{upper(level)} return
$\mu_h \mp z\sqrt{v_h}$ with $z=\Phi^{-1}((1+\text{level})/2)$, computed by
\code{nabia::stats::normal\_quantile}. These are exact intervals under
\eqref{es:eq:ets} with known parameters. They ignore the estimation error in
$\alpha$ and $\sigma^2$, which makes them slightly too narrow in small samples.

\paragraph{Complexity and validation.}
Fitting is a single $O(T)$ pass with $O(T)$ memory (the fitted values), and
forecasting is $O(H)$. The constructor rejects $\alpha\notin(0,1]$, and
\code{fit} rejects an empty or non-finite series. \code{forecast} before
\code{fit} throws \code{std::logic\_error}.

\paragraph{Correspondence with the code.}
\begin{center}
\begin{tabular}{@{}ll@{}}
\toprule
Mathematics & \code{SimpleExponentialSmoothing} \\
\midrule
$\alpha$ & \code{parameters().alpha} \\
$\ell_{T-1}$ & private \code{level\_} \\
$(\ell_{t-1})_{t}$ & \code{fitted\_values()} \\
$\mathrm{SSE}$ & \code{sse()} \\
$(\mu_h, v_h)$ & \code{forecast(H).mean}, \code{.variance} \\
\bottomrule
\end{tabular}
\end{center}

\paragraph{Tests.}
\code{tests/test\_exponential\_smoothing.cpp} covers several cases: a hand-computed series (levels
$10\to15\to22.5$, $\mathrm{SSE}=325$, variances $325/3$ and $1.25\cdot325/3$),
$\alpha=1$ as the naive forecast, invalid $\alpha$, and forecasting before
fitting. The Kalman filter equivalence of Proposition~\ref{es:prop:muth} is tested
in \code{tests/test\_forecasters.cpp}.

%==============================================================================
